\documentclass[11pt]{article}
\usepackage[T1]{fontenc}
\usepackage{lmodern}
\usepackage[margin=0.88in]{geometry}
\usepackage{amsmath,amssymb,amsthm,mathtools,booktabs,array,longtable}
\usepackage{microtype,xcolor,enumitem,fancyhdr}
\usepackage[breaklinks,colorlinks,linkcolor=blue!45!black,urlcolor=blue!45!black,citecolor=blue!45!black]{hyperref}
\usepackage{listings}
\definecolor{navy}{RGB}{23,44,67}
\definecolor{accent}{RGB}{22,100,119}
\lstset{basicstyle=\ttfamily\small,breaklines=true,columns=fullflexible,keepspaces=true,frame=single,rulecolor=\color{black!15}}
\setlength{\parindent}{0pt}
\setlength{\parskip}{5pt}
\setlength{\headheight}{14pt}
\pagestyle{fancy}
\fancyhf{}
\fancyhead[L]{\small\color{navy}Diophantine consistency: descent and proof budgets}
\fancyhead[R]{\small\texttt{2605669a5323}}
\fancyfoot[C]{\thepage}
\setlist{nosep,leftmargin=1.5em}
\newtheorem{theorem}{Theorem}[section]
\newtheorem{lemma}[theorem]{Lemma}
\newtheorem{proposition}[theorem]{Proposition}
\newtheorem{corollary}[theorem]{Corollary}
\theoremstyle{definition}
\newtheorem{definition}[theorem]{Definition}
\theoremstyle{remark}
\newtheorem{remark}[theorem]{Remark}
\newcommand{\Z}{\mathbb Z}
\newcommand{\N}{\mathbb N}
\newcommand{\Q}{\mathbb Q}
\newcommand{\Zp}{\mathbb Z_{>0}}
\newcommand{\Con}{\operatorname{Con}}
\newcommand{\ZFC}{\mathsf{ZFC}}
\newcommand{\Pos}{\operatorname{Pos}}
\newcommand{\bits}{\operatorname{bits}}
\newcommand{\repo}{\texttt{VladimirReshetnikov/ProveIt}}
\newcommand{\pin}{\texttt{2605669a532342f153b13eb6b4a0af4b54f5fa37}}
\newcommand{\file}[1]{\texttt{\detokenize{#1}}}
\allowdisplaybreaks[2]

\begin{document}
\thispagestyle{empty}
{\color{accent}\sffamily\bfseries RESEARCH ITERATION \hfill 23 OCTOBER 2026}
\vspace{1.0cm}
{\color{navy}\Huge\bfseries Diophantine consistency\\[4pt] without redundant\\[4pt] denominator guards\par}
\vspace{0.5cm}
{\Large Pell descent, ordinary strong norms,\\a 31-integer sextic, and radix-budget proof certificates\par}
\vspace{0.5cm}
Prepared for Vladimir Reshetnikov\\
AI-assisted mathematical research; exact arithmetic artifacts included.
\vspace{0.6cm}

\begin{abstract}
We give new explicit arithmetic reductions of the repository's consistency
equations. The eleven-positive-variable construction drops from degree 325
and 113 arithmetic operations to degree \textbf{277} and \textbf{111}
operations. A more substantial rearrangement of the proof yields a
twelve-positive-variable equation of degree \textbf{70} in \textbf{100}
operations, improving the previous degree-82 construction. A compatible
composition of two recent signed-domain arguments yields a sextic in
\textbf{31 unrestricted integers}, improving the 32-integer sextic.
A separate nine-positive-variable construction lowers degree to
\textbf{1,158,240}, at the cost of 811 operations.

The most different research direction is a direct proof-certificate
encoding: a descending coefficient budget replaces an entire coprimality
block by uniqueness of positional notation. Its finite-certificate theorem
is proved, but its final ordinary Diophantine compiler output has not been
materialized or counted. It is not included among the completed polynomial
size records.

The completed equations use exactly specified fixed coefficients and compose
with the named ordinary native-encoding and refutation-calibration proofs
already in ProveIt. We supply full proofs of the new reductions and precise
source boundaries. We do not claim global minimality, a world-first result,
a proof of consistency, or a new end-to-end Lean/Rocq derivation.
\end{abstract}

\vfill
\small
Reviewed snapshot: \pin.\\
The commit time is 24 October in UTC and 23 October in Washington State.\\
Primary deliverable: a self-contained ZIP containing this report, its complete
LaTeX source, literal polynomial graphs, actual coefficient data, generators,
proof notes, source ledgers, and verification scripts.
\newpage
\tableofcontents
\newpage

\section{What is established, and what is being optimized}

\subsection{Logical target}

Fix the recursive syntactic ZFC refutation predicate used by the inherited
repository calibration. Write $\Con(\ZFC)$ for the arithmetical assertion
that no finite refutation exists. The target is a fixed integer polynomial
$P$ and a specified domain $\mathcal D$ for which
\begin{equation}\label{eq:target}
 \ZFC\vdash
 \left[\neg\exists\mathbf x\in\mathcal D^m\;P(\mathbf x)=0\right]
 \longleftrightarrow\Con(\ZFC).
\end{equation}
The two sides of this equivalence are not asserted separately. The ordinary
proof proceeds by constructing finite refutations from zeros and constructing
zeros from finite refutations. It does not assume that ZFC is consistent or
arithmetically sound.

The new arguments below are elementary finite arithmetic and Pell arguments
formalizable in ZFC. Their endpoint uses the repository's already supplied
ordinary compiler/calibration theorems. The repository is separately working
toward a kernel-checked derivation in its encoded object-language ZFC; that
formalization remains distinct from an ordinary mathematical proof. This
archive does not close that formalization task.

\subsection{A measured improvement has several coordinates}

The pinned proposal review \cite{review} distinguishes variable count,
domain, total degree, arithmetic-circuit size, and fixed-coefficient size.
That distinction matters here. The following are comparisons with specific
frontier points, not an ordering of every equation by a single notion of
``simplest.''

\begin{center}
\small
\begin{tabular}{@{}lrrr rrr@{}}
\toprule
 &\multicolumn{3}{c}{Pinned predecessor}&\multicolumn{3}{c}{This contribution}\\
Domain &Variables&Degree&Ops.&Variables&Degree&Ops.\\
\midrule
Positive integers &11&325&113&11&277&111\\
Positive integers &12&82&102&12&70&100\\
Integers, sextic &32&6&--&31&6&--\\
Positive integers &9&1,379,730&478&9&1,158,240&811\\
\bottomrule
\end{tabular}
\end{center}

The eleven-variable degree falls by $48/325\approx14.77\%$. The
twelve-variable degree falls by $12/82\approx14.63\%$, with two fewer
operations. The nine-variable degree falls by approximately $16.05\%$,
while its operation count increases; it is a degree/size tradeoff. The
repository's 443-operation nine-positive-variable target therefore remains
the better choice when operation count is the priority.

One operation means a live binary addition, subtraction, or multiplication;
squaring and scalar multiplication are charged. Fixed numerals and reuse of
a computed register are free. These counts do not measure bit complexity.
Named registers are polynomial abbreviations, not existential variables.
The integer sextic has 18 cubic-or-lower residuals and 815 terms when the
fixed coefficients are treated as coefficient symbols; no optimized scalar
circuit count is claimed for it.

\subsection{Contributions and inherited work}

The rational integrality lemma comes from package 69. The three-root norm
comes from package 70. The adaptive carrier comes from package 73. The
ordinary strong-rank lemma predates those packages. The present positive
equations exploit their compatibility and change the order in which facts
are established. Similarly, the integer sextic combines package 75's signed
bootstrap with package 71's determinant guard.

The direct coefficient-budget certificate is a different architecture
relative to the inspected direct source in package 68. Its basic mathematical
tools---coefficient bounds, positional uniqueness, and factorization---are
classical. The new contribution is their particular use to remove the source's
entire common-multiple block. We make no claim that these general tools are new.

\section{Two elementary lemmas that remove expensive guards}\label{sec:lemmas}

\subsection{Rational integrality by descent}

\begin{lemma}[Adjacent-coefficient Pell descent; package 69]\label{lem:descent}
Let $Q>1$ and $V$ be integers, and let $y\in\Q$. If
\begin{equation}\label{eq:adjacent}
 QV^2-(Q-1)y^2=1,
\end{equation}
then $y\in\Z$.
\end{lemma}
\begin{proof}
The equation excludes $V=0$ and $y=0$. Take absolute values and write
$y=a/m$ in lowest terms, $m>0$. Since $(Q-1)a^2/m^2$ is an integer,
$m^2\mid Q-1$. If $V=1$, then $y=1$ and $m=1$.

For $V>1$ set
\[
 V'=(2Q-1)V-2(Q-1)y,\qquad
 y'=(2Q-1)y-2QV.
\]
The first coordinate is an integer because $m\mid Q-1$. Direct expansion
preserves \eqref{eq:adjacent}. The inequalities $y>V$ and
$QV>(Q-1)y$ imply $-V<V'<V$; the new norm excludes $V'=0$.
Moreover $2Q-1\equiv1\pmod m$, so $y'$ has the same reduced denominator
$m$. Strong induction on the positive integer $|V|$ reaches $V=1$ without
changing this denominator. Thus $m=1$.
\end{proof}

The adjacency and the unit right side are essential. For example,
$2\cdot1^2-4(1/2)^2=1$ has a nonintegral second coordinate. A generic
resultant is insufficient: it initially recovers rational witnesses, and
the descent is what promotes them to integers.

\subsection{A rational sum of three square roots}

\begin{lemma}\label{lem:three}
Let $d_1,d_2,d_3$ be positive rational numbers, pairwise distinct, with no
zero signed sum of their positive square roots. If a signed sum of those
roots is rational, then each signed root is rational.
\end{lemma}
\begin{proof}
Write the signed roots as $x_1,x_2,x_3$, their squares as $a,b,c$, and
their rational nonzero sum as $s$. Set
\[
 u=\frac{s^2-a-b-c}{2},\qquad
 v=\frac{u^2-ab-ac-bc}{2s}.
\]
Then $u=x_1x_2+x_1x_3+x_2x_3$ and $v=x_1x_2x_3$. Consequently
\[
 x_1=\frac{a(a+u)}{as+v},\qquad
 as+v=x_1(x_1+x_2)(x_1+x_3).
\]
The denominator is nonzero: cancellation with another root would make
their squares equal. The two cyclic formulas prove the other cases.
\end{proof}

The universal conjugate identity used below is
\begin{align}\label{eq:octic}
 F_8(x;a,b,c)
 &=\bigl[(x^2-a-b-c)^2-4(ab+ac+bc)\bigr]^2-64abcx^2\\
 &=\prod_{\varepsilon\in\{-1,1\}^3}
   (x-\varepsilon_1\sqrt a-\varepsilon_2\sqrt b-\varepsilon_3\sqrt c).
 \nonumber
\end{align}
The supplied verifier proves this identity by exact polynomial multiplication.

\section{An eleven-positive-variable polynomial of degree 277}\label{sec:eleven}

\subsection{Fixed coefficients and common registers}

Use the complete compact44 coefficient tuple
$(\beta,K_0,M_C,M_F)$ of package 44 \cite{native44}, with
\[
 \beta=2^{590069}-1,\qquad M_F=M_{F,\mathrm{native}}+\beta.
\]
The high mask is shifted exactly once. Appendix~\ref{sec:coefficients}
specifies the tuple by finite arithmetic and a literal action seed; actual
integer payloads and deterministic generators are also included. The four
entries must be used together.

The eleven independent variables, all in $\Zp$, are
\[
 j,F,\alpha,i,T,X,Y,n,\eta,\zeta,\gamma.
\]
Define the following polynomial registers:
\begin{align*}
 J&=3j+1,& q&=\beta J+1,& C&=q-F-\alpha,\\
 E&=XY,& k&=\eta+\zeta,& c&=kY+\eta,\\
 a&=E+Y,& H&=4a+3,& \Delta&=a^2+H=(a+2)^2-1,\\
 \ell&=EkY,& I&=ic^2,& D_f&=1+\Delta I^2,\\
 Q&=\Delta^2 I^2,& \delta&=Q-1,& D_\tau&=\ell(\ell+k)+1.
\end{align*}
Also set
\begin{align*}
 R&=((q-1)(q-F)+\alpha)(q^2-1)+(M_C+qM_F)J,\\
 \widehat R&=q^3(q+F+\alpha),& M&=q^3(q-1),\\
 B_t&=(qK_0+X)C-qF,& V&=D_f(cT-R)-c,\\
 L&=(D_f+1)(cT+\widehat R),& W&=1+(q-1)L,\\
 Z&=M(\ell+k+n+L)+(q-1)Y-q^3WB_t.
\end{align*}
The adaptive main-root residual is
\begin{align*}
 r_0&=k-R-1,& S&=4qY,& U&=qS,\\
 D_*&=U(X+ac)+H(U\gamma+SX-r_0),\\
 A_1&=D_*^2-U^2(1+\Delta c^2).
\end{align*}

For the norm write
\begin{align*}
 A&=M^2D_\tau,& B_f&=M^2D_f,& C_n&=M^2(QV^2-1),\\
 N_c&=\delta^2\bigl[(Z^2-A-B_f)^2-4AB_f\bigr]
       -2\delta C_n(Z^2+A+B_f)+C_n^2,\\
 N_8&=N_c^2-64\delta^3AB_fC_nZ^2.
\end{align*}
Our equation is the single, division-free polynomial equation
\begin{equation}\boxed{P_{11}=N_8+A_1+N_8A_1=0.}\label{eq:p11}\end{equation}
Its actual graph is \file{positive/p11_positive_degree277.json}.

\subsection{Unconditional root separation and the smaller majorant}

Before imposing any residual, put
\[
 t=\sqrt{D_\tau},\quad f=\sqrt{D_f},\quad
 D_y=\frac{QV^2-1}{Q-1},\quad y=\sqrt{D_y}.
\]
These are proof notation only. Positive inputs give $c\ge3$, $a\ge2$,
$\Delta\ge15$, and $I\ge c^2$. In particular
\[
 1<D_\tau<D_f<Q,\qquad
 t<\ell+k<(a+1)c,\qquad f>(a+1)c^2.
\]
Let $b=cT-R$. If $b=0$, then $V=-c$; otherwise
$|V|\ge D_f-c>1$. Thus
\[
 D_y=V^2+\frac{V^2-1}{Q-1}>V^2>1.
\]
If $D_y$ equaled either $D_\tau$ or $D_f$, the equation
$Q(D_y-V^2)=D_y-1$ would put a nonzero integer multiple of $Q$ strictly
between $0$ and $Q$. Hence the three radicands are pairwise distinct.

If $b=0$, then $y<c+1<f-t$; if $b\ne0$, then
$y>|V|\ge f^2-c>2f>f+t$. One root therefore strictly dominates the other
two, so no signed sum is zero. Lemma~\ref{lem:three} applies whenever the
octic norm vanishes.

The authentic masks satisfy $0<M_C<\beta$ and $0<M_F<2\beta$. As in
package 70, these imply $|R|<\widehat R$ for every positive input tuple,
without assuming $C>0$. Put $B_*=cT+\widehat R$ and
$\overline V=D_fB_*$. Then
\[
 |V|\le D_f(cT+|R|)+c,\qquad
 \overline V-|V|\ge D_f-c>f.
\]
Also $Q-1>D_f$ and
\[
 y-|V|=\frac{V^2-1}{(Q-1)(y+|V|)}
       <\frac{|V|}{Q-1}<B_*.
\]
Consequently
\begin{equation}\label{eq:sharp-majorant}
 f+y<L=(D_f+1)B_*,\qquad L>k.
\end{equation}
This sharpens the older $3D_fB_*$ majorant without adding an operation.

\subsection{Soundness: why the deleted factor is redundant}

Both $N_8$ and $A_1$ are $0$ or $1$ modulo $4$. Therefore
$(1+N_8)(1+A_1)=1$ forces both factors to be $1$: a factor $-1$ would
make its residual $-2\equiv2\pmod4$. Thus \eqref{eq:p11} implies
$N_8=A_1=0$.

By \eqref{eq:octic}, the zero norm gives signed roots with
\[
 Z/M=\tau_*+f_*+y_*,\qquad
 \tau_*^2=D_\tau,\quad f_*^2=D_f,\quad y_*^2=D_y.
\]
Lemma~\ref{lem:three} makes all three roots rational. The first two are
integers, since their squares are integers. For the third,
\[
 QV^2-(Q-1)y_*^2=1
\]
has integral $Q,V$, so Lemma~\ref{lem:descent} gives $y_*\in\Z$.

This is the precise improvement. Package 70 inserted $q(q-1)$ into $I$
to recover $y_*$ by a denominator-coprimality argument. Descent proves its
integrality directly, so those two multiplications and their degree cost
are unnecessary.

Set
\[
 f_{\rm old}=|f_*|,\quad u=|y_*|-V,\quad
 r=n+L-f_*-y_*.
\]
The inequalities above give $u,r>0$. The signed roots, not their absolute
values, must be used in $r$. Subtracting $M(f_*+y_*)$ from $Z$ yields
\[
 Z_0=M(\ell+k+r)+(q-1)Y-q^3WB_t=M\tau_*.
\]
We have recovered the weighted first carrier, the normalized strong norm,
and the auxiliary norm:
\begin{align*}
 Z_0^2&=M^2D_\tau,\\
 f_{\rm old}^2-\Delta(ic^2)^2&=1,\\
 (V+u)^2-\Delta(f_{\rm old}^2-1)u(2V+u)&=1.
\end{align*}
Together with $A_1=0$, these are precisely the weighted/adaptive source
interface proved in package 70 \cite{source70}. Its soundness is valid for
every positive integral weight; here that weight is $L$. It yields a finite
ZFC refutation through the inherited compact44 native interface.

\subsection{Completeness and order of choices}

A finite refutation supplies canonical normalized source witnesses
\cite{source61,native44}. Retain their $q,X,Y,k,c,R,f,i$ and choose the
positive auxiliary pair as in the source. In particular, $c\mid I$ permits
the usual odd-index auxiliary construction. Compute $L,W$ only after these
choices. No new divisibility by $q(q-1)$ is needed.

Write $s=Y/q^3$, and let $t_0=B_t/(q-1)=q(z_{\rm source}-1)\ge q$.
The first positive root satisfies $\ell<\tau<\ell+k$, and $0<s<k$.
Choose
\[
 r=Wt_0+\tau-\ell-k-s,\qquad
 n=r-L+f+y.
\]
Since $q\ge16$ and $L>k$, we have
$r>qW-2k>L$, so $n>0$. The weighted carrier and adaptive main carrier
lift as in the source, and $Z=M(\tau+f+y)$, giving $N_8=A_1=0$.
This establishes both existence implications without assuming consistency.

\subsection{Exact degree and size}

The relevant degrees are
\[
\begin{array}{c|rrrrrrrrrr}
\text{register}&I&D_f&Q&V&L&W&Z&N_c&N_8&A_1\\\hline
\deg&5&14&18&18&18&19&24&132&264&13.
\end{array}
\]
The apparent degree-14 leading terms in $A_1$ cancel; exact expansion gives
degree 13. Let $q_0=3\beta j$, $k_0=\eta+\zeta$, and
$B_0=(K_0q_0+X)(q_0-F-\alpha)-q_0F$. Then
\begin{align*}
 [Q]_{18}&=i^2k_0^4X^4Y^8,\\
 [Z]_{24}&=-i^2q_0^7k_0^4X^2Y^6(q_0+F+\alpha)B_0,\\
 [A_1]_{13}&=32q_0^2k_0X^2Y^4
    \bigl(4q_0^2Y\gamma+q_0^3(q_0-F)\bigr).
\end{align*}
The leading form of $N_8$ is $[Q]_{18}^4[Z]_{24}^8$. In $P_{11}$ the
monomial
\[
 j^{60}i^{24}\eta^{49}X^{42}Y^{85}\alpha^{16}\gamma
\]
has coefficient $2^7(3\beta)^{60}>0$. Its total degree is 277, independently
of the other three fixed ports. Thus the degree is exactly 277 on the actual
coefficient slice. The delivered graph has 64 multiplications, 30 additions,
and 17 subtractions: 111 operations, or 108 before the finalizer.

\section{A twelve-positive-variable polynomial of degree 70}\label{sec:twelve}

\subsection{The ordinary strong norm becomes available earlier}

The normalized equation $f^2-\Delta(ic^2)^2=1$ is convenient because it
classifies $f$ immediately. The older ordinary equation
\[
 (ic^2)^2=\Delta(f^2-1)
\]
has lower degree, but its rank-recovery proof needs the additional estimate
$c>(a+2)\Delta^2$. Applying that estimate only after auxiliary index
recovery would be circular. The adaptive carrier yields an independent
lower bound on the first Pell index, so the estimate can be proved before
using the strong equation. This changes the dependency order and licenses
the smaller norm.

\subsection{Complete equation}

Use the same compact44 coefficients. The twelve positive variables are
\[
 j,F,\alpha,t,T,X,Y,r,\eta,\zeta,n,\gamma.
\]
Set $J=6j+1$, $q=\beta J+1$, and retain $C,E,k,c,a,H,\Delta,\ell,R,B_t$
from Section~\ref{sec:eleven}. This time define
\begin{align*}
 U&=q^2Y,&v&=k-R-1,&W&=1+(q-1)k,&M&=q^3(q-1),\\
 D_*&=U(X+ac)+H(U\gamma+qE-v),\\
 Z&=M(\ell+k+r)+(q-1)Y-q^3WB_t,\\
 A_0&=Z^2-M^2[\ell(\ell+k)+1],\\
 A_1&=D_*^2-U^2(1+\Delta c^2).
\end{align*}
Define the ordinary norm merger by
\begin{align*}
 f&=t+1,&\rho&=f^2-1,&Q&=\Delta\rho,&d&=c^4,\\
 V&=f^2(cT-R)-c,\\
 h(z)&=(1-Q)z^2+2(1-Q)Vz+V^2-1,\\
 b_N&=2(1-Q)(n+V),& A_N&=d\,h(n)+(1-Q)Q,\\
 N&=A_N^2-d\,b_N^2Q.
\end{align*}
The equation is
\begin{equation}\boxed{P_{12}=(1+A_0)(1+A_1)(1+N)-1=0.}\label{eq:p12}\end{equation}
Its graph is \file{positive/p12_positive_degree70.json}. Computing $qE$
reuses the existing $E=XY$, saving a multiplication compared with a literal
left-to-right evaluation of $qXY$.

\subsection{Recovering the outer coordinates without the strong norm}

Because $\beta,J$ are odd, $q$ is even. Thus $A_0,A_1$ are squares modulo
4; $b_N$ is even, so $N$ is a square modulo 4 as well. The unit-product
argument gives $A_0=A_1=N=0$.

The first two residuals alone give the following facts
\cite{source61,source73}. We include the order because it is essential.
From $A_0=0$, the rational square root $Z/M$ is integral. Reducing $M\mid Z$
modulo $q^3$ and $q-1$ gives
\[
 q^3\mid Y,\qquad q-1\mid B_t.
\]
Write $s=Y/q^3>0$ and $b=B_t/(q-1)$. The identity
$Z/M=\ell+k+r+s-Wb$, together with $|Z/M|<\ell+k$, forces $b>0$.
It follows that $B_t>0$ and $C>0$. The mask bounds now imply
\[
 q\ge16,\qquad 3q<R<q^4-2.
\]
Since $q\equiv2\pmod3$ and $q^3\mid Y$, we have $\gcd(q,H)=1$.

The estimate $U\ge q^5$ shows $D_*>0$. From $A_1=0$, the rational root
$D_*/U$ is an integer $g>0$. Its divisibility condition gives
$U\mid H(qXY-v)$. Since $q^5\mid U$ and $q^4\mid qXY$, we obtain
\begin{equation}\label{eq:early-congruence}
 q^4\mid v=k-R-1.
\end{equation}

Put $A=a+2$ and $P_0=2XY^2+1$. Positive Pell classification gives
indices $m_1,p\ge1$ with
\[
 |Z/M|=\chi_{P_0}(m_1),\qquad
 k=2\psi_{P_0}(m_1),\qquad c=\psi_A(p).
\]
Here $P_0>A$ and $c>k$, so $p>m_1$. The first recurrence yields
$k\equiv2m_1\pmod{4XY^2}$, and $q^4\mid4XY^2$. Combining this with
\eqref{eq:early-congruence},
\[
 2m_1\equiv R+1\pmod{q^4}.
\]
Because $0<R+1<q^4$ and $2m_1>0$, the multiplier is nonnegative, hence
\begin{equation}\label{eq:index-lower}
 m_1\ge(R+1)/2\ge24,\qquad p>m_1.
\end{equation}
This lower bound has been obtained before imposing the strong or auxiliary
norms. Pell growth now gives
\[
 c>A\Delta^2,\qquad c>2p,\qquad c>2R.
\]
For the first bound, $p\ge6$ already suffices, since
$\psi_A(6)=32A^5-32A^3+6A>A(A^2-1)^2$. For the third, $m_1\ge2$ gives
$c>kY\ge4P_0Y>8XY^3\ge8q^9>2R$.

\subsection{The merger recovers ordinary integral witnesses}

Before using $N=0$, $f=t+1\ge2$, $Q>1$, and the preceding estimates imply
\[
 V=f^2(cT-R)-c>f^2c/2-c\ge c>1.
\]
Thus $b_N\ne0$. Define rational proof coordinates
\[
 i_*=\frac{A_N}{db_N},\qquad u_*=n-i_*,\qquad y_*=V+u_*.
\]
The zero norm gives $d i_*^2=Q$, and expanding $h(n-i_*)$ gives
$h(u_*)=0$. Equivalently,
\[
 QV^2-(Q-1)y_*^2=1.
\]
Lemma~\ref{lem:descent} gives $y_*\in\Z$. Therefore $u_*\in\Z$ and
$i_*\in\Z$. Set $i=|i_*|>0$ and $y=|y_*|$. We recover
\begin{equation}\label{eq:ordinary-recovered}
 (ic^2)^2=\Delta(f^2-1),\qquad
 (ic^2)^2V^2-((ic^2)^2-1)y^2=1.
\end{equation}
The exact algebraic identity behind completeness is
\begin{equation}\label{eq:ordinary-factorization}
 N\big|_{Q=di^2}=d^2h(n-i)h(n+i).
\end{equation}

\subsection{Ordinary strong-rank recovery and the source lift}

The ordinary strong-rank lemma \cite{ordinary} applies to
\eqref{eq:ordinary-recovered}, $c=\psi_A(p)$, and $c>A\Delta^2$.
It yields
\[
 f=\chi_A(m),\qquad c\mid m,\quad m>2p,
 \qquad ic^2=\Delta\psi_A(m),\qquad c^2\mid f^2-1.
\]
For completeness, the mechanism is as follows. Write
$A=\chi_{F_0}(e)$ for the fundamental Pell parameter and
$\Delta=(F_0^2-1)S^2$, with $S=\psi_{F_0}(e)<A$. The supplied ordinary
strong equation expresses $f=\chi_{F_0}(u)$ and
$ic^2=(F_0^2-1)S\psi_{F_0}(u)$. Strong divisibility, with
$b=\psi_{F_0}(\gcd(u,ep))$, gives $b\ge c/\Delta$. If the gcd were a
proper divisor of $ep$, then
\[
 2b^2<\psi_{F_0}(2\gcd(u,ep))\le Sc<Ac,
\]
contradicting $c>A\Delta^2$. Hence $u=em$ and $p\mid m$. Writing $m=pt$,
the recurrence gives $c\mid\Delta t$. Also
$\gcd(c,\Delta)=\gcd(p,\Delta)$, so $c\mid m$. These facts give the
displayed conclusions. The full inherited statement is retained verbatim.

With $M_*=ic^2$, the signed auxiliary-index lemma
\cite{source61} now applies to
\[
 M_*^2V^2-(M_*^2-1)y^2=1,
 \qquad V\equiv-c\pmod f,\quad V\equiv-R\pmod c.
\]
It yields $p=R$. Returning to \eqref{eq:early-congruence} and using
$0<m_1<p=R<q^4$, the congruence forces $2m_1=R+1$ exactly.
Consequently $v=k-2m_1$ is a positive multiple of $4XY^2$.
Since $U=q^2Y$ divides $4XY^2$, the main-carrier divisibility reduces to
$q\mid HX$, hence $q\mid X$.

Put $w=X/q>0$, $h=v/E>0$, and
$\gamma_{\rm old}=\gamma+w-v/U$. These are integers; Pell growth gives
$\gamma_{\rm old}>0$ and
$g=X+ac+H\gamma_{\rm old}$. The cleared transport has
$B_t=q[(K_0+w)C-F]>0$ and $q-1\mid B_t$, so
\[
 z=\frac{(K_0+w)C-F}{q-1}>0
\]
is an integer. Finally choose $T_{\rm old}=fT$, preserving
$V=c(T_{\rm old}f-1)-Rf^2$.

These coordinates satisfy the \emph{ordinary} source kernel
\cite{ordinary}, which composes with compact44's proved native interface.
One must not claim that dividing $i$ by $\Delta$ always produces a normalized
source witness: that quotient need not be integral. The ordinary source
lift is the correct endpoint.

\subsection{Completeness and exact degree}

A finite refutation gives canonical normalized source witnesses. Replace
its normalized multiplier by
$i_{\rm ordinary}=\Delta i_{\rm normalized}$. This preserves the strong
root and the auxiliary parameter. Choose the positive auxiliary pair with
$u=y-V>0$ and set $n=i_{\rm ordinary}+u>0$ and $t=f-1>0$.
The canonical odd cell count makes $J\equiv1\pmod6$, so the supplied
$j=(J-1)/6$ is positive. The adaptive carrier lifts by
$\gamma=\gamma_{\rm old}-w+v/U>0$, and the first carrier lifts as in
package 73. Identity~\eqref{eq:ordinary-factorization} gives $N=0$.

The exact residual degrees are $(15,13,42)$; the product has degree 70.
Writing $q_0=6\beta j$, $k_0=\eta+\zeta$,
$R_0=q_0^3(q_0-F)$ and
$B_0=(K_0q_0+X)(q_0-F-\alpha)-q_0F$, their leading forms are
\begin{align*}
 [A_0]_{15}&=-2q_0^8k_0^2XY^2B_0,\\
 [A_1]_{13}&=8q_0^2k_0X^2Y^4(q_0^2Y\gamma+R_0),\\
 [N]_{42}&=4k_0^8X^4Y^{12}t^8R_0^2n^2.
\end{align*}
Thus the degree-70 monomial
\[
 j^{18}F^2\eta^{11}X^8Y^{19}t^8n^2\gamma\alpha
\]
has coefficient $64(6\beta)^{18}>0$. The literal graph uses 55
multiplications, 26 additions, and 19 subtractions: 100 operations, or
94 for the three-residual system. Both the upper degree and this nonzero
coefficient are checked by exact sparse polynomial arithmetic.

\input{../sextic31/section.tex}

\input{nine_section.tex}

\input{radix_budget_section.tex}

\section{Reproducibility, proof boundaries, and integration}

\subsection{Contents and primary commands}

The archive is designed to extract as a new research contribution. It does
not modify the upstream repository. Run its standard-library verifier from
the extracted root:
\begin{lstlisting}
python3 verify_all.py
python3 -O verify_all.py
\end{lstlisting}
The positive graphs have independent exact formula checks, conjugate and
resultant identities, degree cancellations, leading-coefficient certificates,
and literal coefficient checks. The sextic verifier reconstructs each
residual from retained source rows and verifies all six eliminations. The
nine-variable verifier checks the polynomial tables and the positivity
certificate used by the new majorant. The direct-certificate verifier
includes positive cases and a counterfeit accepted when its budget is omitted.

The coefficient builders reconstruct actual integers from literal finite
machine data; the integer payloads are also included. Mathematical-byte
hashes and source Git-blob hashes are recorded for reproducibility. A hash
is a checksum, not the mathematical definition of an otherwise missing
coefficient.

\subsection{What the checks do and do not establish}

There are three separate layers of evidence:
\begin{enumerate}
\item Exact finite identities and data: the delivered polynomials, degrees,
operation counts, coefficient payloads, and algebraic transformations.
\item Universal ordinary mathematics: the descent, rank arguments, sign
restoration, canonical witness constructions, and coefficient-budget theorem.
\item The inherited native encoding and syntactic refutation calibration,
followed by the separate project of an end-to-end kernel formalization.
\end{enumerate}
Finite enumeration is evidence for a finite program or an algebraic identity
only within its actual scope. The universal arguments are proved in the
text and the retained source proofs. None of the scripts searches for a
ZFC contradiction or establishes nonhalting by simulation.

The four completed equations in the comparison table are specific
constructions. We have not proved any lower bound that would make one
globally simplest. In particular, a reduction in the degree of a nine-variable
polynomial does not improve the nine-variable count itself.

\subsection{Recommended next mathematical step}

For a compact completed equation, the degree-70 twelve-positive form is the
strongest combined degree-and-operation improvement in this package. For an
unrestricted integer equation with modest degree, the 31-variable sextic is
the useful new target. The nine-variable construction is an alternative when
degree receives more weight than evaluator size.

For a more fundamental change of encoding, compile the direct descending
budget source against the repository's now-complete literal ZFC refutation
checker. Its next gate is precise: materialize and count the ordinary
polynomial obtained from the one bounded-universal block, and transport
its acceptance theorem through the actual encoded ZFC syntax. The source
improvement is already proved, but a final size improvement from that route
must be measured rather than inferred from its source counts.

\appendix
\section{A finite specification of the compact44 coefficients}\label{sec:coefficients}

This appendix supplies a finite arithmetic definition of every fixed
coefficient in $P_{11}$ and $P_{12}$. It is the exact package-44 recipe,
not a new coefficient optimization. The compact54 tuple used by a different
repository target is not substituted here.

\subsection{Literal action seed}

Let $s$ be the following decimal integer, formed by concatenating the printed
lines without spaces. Its length is 1,294 decimal digits and 4,298 bits.
\input{action_seed.tex}
For $0\le q<215$ and $b\in\{0,1\}$, set
\[
 d_{q,b}=\left\lfloor\frac{s}{1024^{2q+b}}\right\rfloor\bmod1024,
 \quad t=d_{q,b}\bmod256,
 \quad v=\left\lfloor d_{q,b}/256\right\rfloor\bmod2.
\]
The move is right if $\lfloor d_{q,b}/512\rfloor=1$ and left otherwise.
These are the 430 literal actions of the package-36 labeling; the initial
state is 93 and the halt state is 215. This is not the repository's separate
dense-state labeling. The delivered verifier reconstructs the original
transition file from these digits.

\subsection{Ordered tile language}

Let $\mathcal R$ and $\mathcal L$ be the sorted sets of targets of right and
left actions, respectively; their cardinalities are 199 and 45. Let
$r(t)=6+\operatorname{rank}_{\mathcal R}(t)$ and
$l(t)=6+199+\operatorname{rank}_{\mathcal L}(t)$, with ranks starting at
zero. Put $h(q,b)=2(q+1)+b$. Form a set of quadruples from the following
finite families:
\begin{align*}
 &(0,h(q,b),r(t),v) &&\text{for right actions},\\
 &(l(t),h(q,b),0,v) &&\text{for left actions},\\
 &(0,h(215,b),0,h(215,b)),\ (0,b,0,b) &&(b=0,1),\\
 &(r(t),b,0,h(t,b)) &&(t\in\mathcal R,\ b=0,1),\\
 &(0,b,l(t),h(t,b)) &&(t\in\mathcal L,\ b=0,1),\\
 &(0,c,0,438) &&(c\in\{0,1,h(215,0),h(215,1)\}),\\
 &(2,434,2,0),\ (2,435,3,1),\ (3,436,4,h(93,0)),\ (4,436,4,0),\\
 &(1,438,1,c) &&(c\in\{434,435,436\}),\\
 &(a,437,b,437) &&(a\in\{0,1,4\},\ b\in\{0,1,2\}),\\
 &(5,0,0,0).&&
\end{align*}
Order the 943 distinct tiles with $(2,435,3,1)$ first, $(5,0,0,0)$ second,
and all others in lexicographic order. Denote them $\tau_0,\ldots,\tau_{942}$.

\subsection{Finite sums defining the four integers}

Put $\rho=2^{67}$, $B=\rho^{8807}$, and $\beta=B-1$. Define the permutation
of $0,\ldots,40$ by the concatenated list
\[
 \pi=(0,\ldots,7,34,8,\ldots,16,35,17,\ldots,24,36,
       25,\ldots,33,37,38,39,40),
\]
and set $w_j=2^{10+\pi_j}$. For the four fields use
\[
 (b_s)=(0,9,19,28),\quad (d_s)=(9,10,9,10),\quad
 (o_s)=(256,512,256,512),\quad (e_s)=(1023,1013,1003,983).
\]
Define a finite table $c_e$ by
\[
 c_i=1+2^{51}-\sum_{s=0}^3\sum_{j=0}^{d_s-1}
 w_{b_s+j}\left(\left\lfloor\frac{(\tau_i)_s+o_s}{2^j}\right\rfloor\bmod2\right)
       -\sum_{j=38}^{40}w_j\quad(0\le i\le942),
\]
together with
\[
 \begin{gathered}
 c_{e_s+j}=w_{b_s+j}\quad(0\le s<4,\ 0\le j<d_s),\\
 c_{1095}=w_{38},\ c_{2195}=w_{39},\ c_{3296}=w_{40},\quad
 c_{1063}=c_{1064}=0.
 \end{gathered}
\]
Let $\mathcal E$ be this table's address set, including the two zero entries.
Set
\begin{align*}
 D_C={}&\rho^{30}+\rho^{62}+\rho^{1100}+\rho^{2201}
       +\sum_{e\in\mathcal E}c_e(\rho^{4393-e}+\rho^{5495-e}),\\
 D_R={}&\rho^{42},\qquad K_0=D_C+BD_R,\\
 M_C={}&B-1-\sum_{e\in\mathcal E\setminus\{1\}}\rho^e-2\rho^{1063}.
\end{align*}
Let
\[
 \mathcal T=\{1013,\ldots,1022\}\cup\{1065,\ldots,1073\}
            \cup\{2195,3296\},\qquad \mu=2^{51}-2.
\]
Then
\begin{align*}
 M_{F,\mathrm{native}}={}&4+\mu(\rho^{4393}+\rho^{5495})
    +\sum_{t\in\mathcal T}\rho^t\\
 &+(2^{65}-1)\sum_{e=8792}^{8804}\rho^e+(2^{19}-1)\rho^{8805},\\
 M_F={}&M_{F,\mathrm{native}}+B-1.
\end{align*}
These are fixed finite sums with fixed exponents. They are not additional
unknowns, variable-exponent terms in the equation, or unevaluated compiler
parameters. All four integers are also supplied as minimal little-endian
bytes and verified against the inherited binding. Their total uncompressed
payload is 295,388 bytes, and their gzip payload is 9,445 bytes.

\section{Source ledger and formalization status}

All repository references used for new comparisons are pinned to
\pin. The complete source ledger in the archive records exact original
paths and Git-blob hashes. Original files keep their original authorship
and license notices. The new text and scripts are marked as AI-assisted.

The primary status distinction is explicit in the pinned repository:
ordinary paper proofs of the compact native interface exist, while complete
Lean native calibration and an encoded internal-ZFC consistency biconditional
remain active work. The latest commit completes a literal numeric checker,
not that final internal transport. This contribution does not silently treat
a standard-model checker theorem as the requested internal derivation.

\begin{thebibliography}{99}
\raggedright
\bibitem{review} ProveIt, \emph{Wave-ten proposals and the current targets},
\href{https://github.com/VladimirReshetnikov/ProveIt/blob/2605669a532342f153b13eb6b4a0af4b54f5fa37/Computability/DiophantineConsistency/PROPOSAL_REVIEW_WAVE10.md}{pinned proposal review}.
\bibitem{native44} ProveIt package 44, \emph{Complement-packed Wang943 native interface},
\href{https://github.com/VladimirReshetnikov/ProveIt/blob/2605669a532342f153b13eb6b4a0af4b54f5fa37/Computability/DiophantineConsistency/con-zfc-diophantine-forms/44-b4-sparse/native/PROOF.md}{proof and exact coefficient recipe}.
\bibitem{source61} ProveIt package 61, \emph{Positive congruence carriers and signed auxiliary-index recovery},
\href{https://github.com/VladimirReshetnikov/ProveIt/tree/2605669a532342f153b13eb6b4a0af4b54f5fa37/Computability/DiophantineConsistency/con-zfc-diophantine-forms/61-congruence-carriers/positive}{source proofs}.
\bibitem{source69} ProveIt package 69, \emph{A Pell descent removes the denominator guards},
\href{https://github.com/VladimirReshetnikov/ProveIt/blob/2605669a532342f153b13eb6b4a0af4b54f5fa37/Computability/DiophantineConsistency/con-zfc-diophantine-forms/69-cones-pell/positive/PROOF.md}{degree-82 proof}.
\bibitem{source70} ProveIt package 70, \emph{Three-root transport: an eleven-positive-variable extension},
\href{https://github.com/VladimirReshetnikov/ProveIt/blob/2605669a532342f153b13eb6b4a0af4b54f5fa37/Computability/DiophantineConsistency/con-zfc-diophantine-forms/70-residue-orientation/eleven/PROOF.md}{degree-325 proof}.
\bibitem{source73} ProveIt package 73, \emph{Adaptive carriers and a twice-centered auxiliary norm},
\href{https://github.com/VladimirReshetnikov/ProveIt/blob/2605669a532342f153b13eb6b4a0af4b54f5fa37/Computability/DiophantineConsistency/con-zfc-diophantine-forms/73-guard-reanchoring/positive/PROOF.md}{adaptive-carrier proof}.
\bibitem{ordinary} ProveIt, \emph{Ordinary strong norms and signed witnesses for the blank computation},
\href{https://github.com/VladimirReshetnikov/ProveIt/blob/2605669a532342f153b13eb6b4a0af4b54f5fa37/Computability/DiophantineConsistency/con-zfc-diophantine-forms/35-cancel12/sources/catalogue/25-hybrid12/upstream/ORDINARY_BLANK_PROOF.md}{retained ordinary-kernel proof}.
\bibitem{source71} ProveIt package 71, \emph{Calibrated determinant guards},
\href{https://github.com/VladimirReshetnikov/ProveIt/tree/2605669a532342f153b13eb6b4a0af4b54f5fa37/Computability/DiophantineConsistency/con-zfc-diophantine-forms/71-determinant-guards}{proofs and nine-variable predecessor}.
\bibitem{source75} ProveIt package 75, \emph{Signed bootstrap},
\href{https://github.com/VladimirReshetnikov/ProveIt/tree/2605669a532342f153b13eb6b4a0af4b54f5fa37/Computability/DiophantineConsistency/con-zfc-diophantine-forms/75-signed-bootstrap}{signed-power proof}.
\bibitem{source68} ProveIt package 68, \emph{Commutative affine-product certificates},
\href{https://github.com/VladimirReshetnikov/ProveIt/blob/2605669a532342f153b13eb6b4a0af4b54f5fa37/Computability/DiophantineConsistency/con-zfc-diophantine-forms/68-joint-positivity/direct/COMMUTATIVE_CERTIFICATES.md}{direct source and flow theorem}.
\bibitem{jones} J. P. Jones, \emph{Universal Diophantine equation},
Journal of Symbolic Logic \textbf{47} (1982), 549--571.
\href{https://doi.org/10.2307/2273588}{doi:10.2307/2273588}.
\end{thebibliography}
\end{document}
